\documentclass[10pt]{article}
\usepackage[a4paper,margin=0.78in]{geometry}
\usepackage{amsmath,amssymb,booktabs,microtype}
\usepackage[T1]{fontenc}
\usepackage{lmodern}
\usepackage{enumitem}
\setlist{nosep,leftmargin=*}
\title{A Three-Site Variational H\'enon Ring:\\ An Exact Fourier-Mode Period-Two Witness}
\author{Route-A finite certificate C111}
\date{}
\newcommand{\Q}{\mathbb{Q}}
\begin{document}
\maketitle
\begin{abstract}
We give a small exact certificate for a genuinely three-site variational H\'enon lattice. A triangular coupling graph, with $a=7$ and $\kappa=1/5$, yields a polynomial exact-symplectic map over $\Q$. Two synchronous fixed points and one primitive synchronous period-two orbit are verified by independent rational arithmetic. The two-step Jacobian splits into one longitudinal and two transverse Laplacian modes, giving the factorisation $(1+7z+z^2)(1+106z/25+z^2)^2$ for the finite monodromy polynomial. This is deliberately a low-period witness; no global zeta or arithmetic claim is made.
\end{abstract}
\section{Scope and model}
This paper records package C111 in the exploratory Route-A line. The claim is
an auditable finite calculation, not a global zeta theorem. Let $q,p\in
\Q^3$ and let
\[
 L=\begin{pmatrix}2&-1&-1\\-1&2&-1\\-1&-1&2\end{pmatrix},\qquad
 U(q)=\sum_{i=1}^3\left(\frac{7}{2}q_i^2-\frac13q_i^3\right)
 -\frac1{10}\sum_{\{i,j\}}(q_i-q_j)^2 .
\]
The sum has the three edges of the $3$-cycle. The map is
\begin{equation}
 F(q,p)=(G(q)-p,q),\qquad G(q)=\nabla U(q).
 \label{eq:map}
\end{equation}
All computations use exact fractions. The evidence file and scripts are
listed in the reproducibility paragraph below.

\section{Exact variational identities}
The Hessian and Jacobian are
\[
 H(q)=\operatorname{diag}(7-2q_i)-\tfrac15L,\qquad
 J(q)=\begin{pmatrix}H(q)&-I_3\\ I_3&0\end{pmatrix}.
\]
With $\Omega=\left(\begin{smallmatrix}0&I_3\\-I_3&0\end{smallmatrix}\right)$,
direct multiplication gives $J(q)^\mathsf{T}\Omega J(q)=\Omega$ and
$\det J(q)=1$. The coordinate swap $R(q,p)=(p,q)$ is an involutory
reversor: $RFR=F^{-1}$. For the canonical one-form
$\lambda=q\mathbin{\cdot}dp$ one has the exact primitive identity
\[
 F^*\lambda-\lambda=d\bigl(U(q)-p\mathbin{\cdot}q\bigr).
\]
These are polynomial identities, checked both by the producer/checker and by
the independent SymPy script.

\section{Certified orbit witnesses}
The synchronous subspace is invariant because $L(1,1,1)^\mathsf{T}=0$.
The two fixed states and the period-two states are
\[
\begin{aligned}
 P_0&=((0,0,0),(0,0,0)),\\
 P_5&=((5,5,5),(5,5,5)),\\
 C_3&=((3,3,3),(6,6,6)),\qquad
 C_6=((6,6,6),(3,3,3)).
\end{aligned}
\]
Substitution into \eqref{eq:map} gives $F(P_i)=P_i$ and
$F(C_3)=C_6$, $F(C_6)=C_3$. Since $C_3\ne C_6$, this is a primitive
period-two witness. No completeness statement is made about other fixed or
periodic points.

\section{Longitudinal and transverse modes}
The Laplacian eigenvalues are $0,3,3$. At a synchronous point with scalar
coordinate $s$, the Hessian eigenvalue in mode $\ell$ is
\[
 h_s(\ell)=7-2s-\tfrac15\ell .
\]
Thus the two points of the cycle have mode pairs
\[
 (h_3(0),h_6(0))=(1,-5),\qquad
 (h_3(3),h_6(3))=(2/5,-28/5).
\]
For a scalar mode, the two-step matrix
$\left(\begin{smallmatrix}h_6&-1\\1&0\end{smallmatrix}\right)
\left(\begin{smallmatrix}h_3&-1\\1&0\end{smallmatrix}\right)$ has determinant one
and trace $h_3h_6-2$. Therefore the mode traces are
\[
 t_0=-7,\qquad t_\perp=-\frac{106}{25}
 \quad\text{(multiplicity two)}.
\]
The direct $6\times6$ monodromy and the mode reconstruction agree exactly:
\begin{equation}
 \det(I-zM)=(1-t_0z+z^2)(1-t_\perp z+z^2)^2
 =1+\frac{387}{25}z+\frac{50211}{625}z^2
 +\frac{98002}{625}z^3+\frac{50211}{625}z^4
 +\frac{387}{25}z^5+z^6 .
 \label{eq:poly}
\end{equation}

\begin{table}[t]
\centering
\caption{Exact two-step mode data on the period-two witness.}
\begin{tabular}{@{}lccc@{}}
\toprule
mode & Laplacian eigenvalue & $(h_3,h_6)$ & trace \\
\midrule
longitudinal & $0$ & $(1,-5)$ & $-7$\\
transverse 1 & $3$ & $(2/5,-28/5)$ & $-106/25$\\
transverse 2 & $3$ & $(2/5,-28/5)$ & $-106/25$\\
\bottomrule
\end{tabular}
\end{table}

As a finite coupling control, setting $\kappa=0$ at the same synchronous
states gives the coefficient vector
\[
 (1,21,150,385,150,21,1),
\]
instead of the coefficients in \eqref{eq:poly}; the trace changes by
$138/25$. This comparison isolates a coupling effect without interpreting
either polynomial as a Fredholm determinant.
\section{Reproducibility and limits}
From the package directory, run
\begin{verbatim}
python3 code/c111_three_site_producer.py
python3 code/c111_three_site_checker.py
python3 code/c111_sympy_crosscheck.py
python3 code/c111_replay.py
python3 code/c111_mutation.py
\end{verbatim}
The canonical evidence is
\texttt{results/c111\_three\_site\_evidence.json}; the mutation audit rejects
all twelve semantic edits. The release manifest records hashes for every
source, result, and paper artifact.

The route verdict is $A1=\texttt{A1\_WEAK}$ (certified low-period witnesses),
$A2=\texttt{A2\_FAIL}$ (operator owner open), $A3$ not addressed, and $A4$
failed; overall status is \texttt{ROUTE\_A\_EXPLORATORY}. In particular, this
package does not claim a complete primitive-orbit atlas, a global transfer or
Fredholm operator, analytic continuation, Euler factors, root numbers,
automorphy, or a Hilbert--P\'olya operator. Its scope literal is
\texttt{NO\_BAD\_EULER\_OR\_ROOT\_NUMBER}.
\end{document}
